%%%%%
%
%% Dans ce fichier, les �ventuels accents sont cod�s
%% en ISO-8859-15 (aka latin9)
%% La ligne suivante donne trois exemples 
%% e accent aigu : �, e accent grave : �, e-dans-l'o : �
%%%%
%
\section{Estimation de la qualit� de la mesure du vent relatif}

\itemstyle{Op\'erateur~: } \operfmt{Qualite\_vent\_relatif}
\index{Qualite\_vent\_relatif}

%%%%%OBJET%%%%%
\itemstyle{Objet~: }
 Ce programme effectue un certain nombre de tests de coh�rence
sur les mesures du vent relatif � une plate-forme mobile (typiquement, 
un avion), et indique les r�sultats dans une s�rie de "flags".

%%%%%ARGUMENTS%%%%
\bigskip
%%%%%ENTREES%%%%
\begin{tabarg}
\multicolumn{3}{l}{\titretab{ Entr\'ees}}\\
\hline
$dp_{1}$ & Vecteur, \texttt{float} & Premi�re mesure de pression diff\'erentielle ($hPa$)\\
$dp_{2}$ & Vecteur, \texttt{float} & Seconde mesure de pression diff\'erentielle ($hPa$)\\
$incid\_rdm$ & Vecteur, \texttt{float} & Incidence radome ($radian$)\\
$incid\_gir$ & Vecteur, \texttt{float} & Incidence girouette ($radian$)\\
$dp_i$ & Vecteur, \texttt{float} & Pression diff\'erentielle incidence radome ($hPa$)\\
$dp_j$ & Vecteur, \texttt{float} & Pression diff\'erentielle d\'erapage radome ($hPa$)\\
\end{tabarg}

\medskip
%%%%%SORTIES%%%%
\begin{tabarg}
\multicolumn{3}{l}{\titretab{ Sorties}}\\
\hline
$flag_{vent}$ & Vecteur, \texttt{int\_32} & Indice de qualit\'e de la mesure du vent ()\\
\end{tabarg}



%%%%%FORMULE%%%%
\itemstyle{Formule}

%%%% Qualit� de la mesure du vent relatif

Si $dp_{1} > 8$ alors \newline
\begin{tabular}{rl}
Test �cart vitesses & si $|dp_{1} - dp_{2}| / dp_{1} > 0.03$ alors $flag_1 = 0001$ (alerte)\\
       & si $|dp_{1} - dp_{2}| / dp_{1} > 0.07$ alors $flag_1 = 0002$ (d�tection)\\
Test �cart incidences& si $|incid_{rdm} - incid_{gir}| >3$ alors $flag_2 = 0010$ \\
Test gamme incidence& si $dpi / dp_{1} < -0.51$ ou $dpi / dp_{1} > 0.48$ alors $flag_3 = 0100$ \\
Test gamme d�rapage& si $|dpj / dp_{1}| > 0.4$ alors $flag_4 = 1000$ \\
\end{tabular}

$flag_{vent} = flag_1 + flag_2 + flag_3 + flag_4$


%%%%%AUTEUR%%%%%
\itemstyle{Auteur~:}
CNRM/GMEI/TRAMM

%%%%%REFERENCES%%%%%
\itemstyle{ R\'ef\'erences~:}
